\chapter[Lezione 2 --- 30 settembre 2026]{Lezione 2: numeri interi, razionali e fattoriale}
\label{chap:lezione-02}

\section{Numeri interi}
\label{sec:lez02-numeri-interi}

Dato un insieme \(A\), si indica con \(-A\) l'insieme degli opposti dei
suoi elementi:
\[
  -A=\{-a:a\in A\}.
\]
Poiché \(0\in\N\), l'insieme \(-\N\) contiene \(0\) e gli opposti dei
naturali positivi. L'unione di \(\N\) e \(-\N\) permette quindi di
introdurre l'insieme dei numeri interi.

\begin{definition}[Numeri interi]
\label{def:numeri-interi}
\[
  \Z=\N\cup(-\N).
\]
\end{definition}

Si assume l'\emph{assioma di esistenza e unicità dell'opposto}, secondo cui
ogni intero ammette un unico opposto, anch'esso appartenente a \(\Z\). In
forma più precisa, vale
\[
  \forall m\in\Z\;\exists!\,n\in\Z
  \quad\text{tale che}\quad
  m+n=0.
\]
Tale unico intero è indicato con \(-m\) ed è detto \emph{opposto} di \(m\).

L'ordinamento è compatibile con la moltiplicazione per un intero positivo;
se invece il fattore è negativo, il verso della disuguaglianza si inverte.

\begin{proposition}[Compatibilità dell'ordine in \(\Z\)]
\label{prop:compatibilita-ordine-interi}
Per ogni \(a,b,c\in\Z\), se \(a\le b\), allora
\[
  c>0\Longrightarrow ac\le bc,
  \qquad
  c<0\Longrightarrow ac\ge bc.
\]
\end{proposition}

In particolare, moltiplicare una disuguaglianza per un numero negativo
inverte il verso. Se \(c=0\), invece, si ha
\[
  ac=bc=0.
\]

Si pone allora il problema di stabilire se, dato \(a\in\Z\), esista
\(b\in\Z\) tale che
\[
  ab=1.
\]
Quando un tale elemento esiste, esso è detto \emph{inverso moltiplicativo}
di \(a\) e viene indicato con \(a^{-1}\), ossia
\[
  a^{-1}:=b.
\]
In \(\Z\), tuttavia, l'inverso moltiplicativo non esiste per ogni elemento.

\begin{proposition}[Inversi moltiplicativi in \(\Z\)]
\label{prop:inversi-interi}
Gli unici interi che ammettono un inverso moltiplicativo in \(\Z\) sono
\(-1\) e \(1\). In particolare,
\[
  (-1)^{-1}=-1,
  \qquad
  1^{-1}=1.
\]
\end{proposition}

\begin{proof}
Siano \(a,b\in\Z\) tali che \(ab=1\). Il prodotto è positivo, perciò
\(a\) e \(b\) hanno lo stesso segno. Se sono entrambi positivi, sono
interi almeno uguali a \(1\); affinché il loro prodotto sia \(1\), devono
quindi essere entrambi uguali a \(1\). Se sono entrambi negativi, gli
interi positivi \(-a\) e \(-b\) hanno prodotto \(1\), perciò
\(a=b=-1\). Viceversa, \((-1)(-1)=1\) e \(1\cdot1=1\), da cui seguono
le uguaglianze enunciate.
\end{proof}

\section{Numeri razionali}
\label{sec:lez02-numeri-razionali}

Come mostrato nella \cref{prop:inversi-interi}, in \(\Z\) soltanto
\(-1\) e \(1\) ammettono un inverso moltiplicativo. Ciò motiva
l'introduzione di un insieme numerico più ampio, nel quale sia possibile
rappresentare anche rapporti tra numeri interi.

Nella definizione dei numeri razionali sarà utile escludere lo zero
dall'insieme dei possibili denominatori. A questo scopo, per due insiemi
\(A\) e \(B\), si introduce la differenza insiemistica
\[
  A\setminus B=\{x\in A:x\notin B\},
\]
ossia l'insieme degli elementi di \(A\) che non appartengono a \(B\).

In particolare,
\[
  \N\setminus\{0\}
\]
indica l'insieme dei numeri naturali diversi da zero.

\begin{definition}[Numeri razionali]
\label{def:numeri-razionali}
L'insieme dei numeri razionali è
\[
  \Q=\left\{\frac{m}{n}\,\middle|\,
    m\in\Z,\ n\in\N\setminus\{0\}\right\}.
\]
\end{definition}

L'introduzione dei numeri razionali estende quindi gli insiemi numerici
considerati finora, secondo la catena di inclusioni
\[
  \N\subseteq\Z\subseteq\Q.
\]

Il denominatore viene scelto positivo; un eventuale segno negativo si
riporta quindi al numeratore. Uno stesso numero razionale può essere
rappresentato da frazioni diverse, come \(\frac12=\frac24\). Riducendo
la frazione ai minimi termini si ottiene una rappresentazione con
\(\operatorname{mcd}(m,n)=1\), dove il massimo comune divisore è inteso
come quello positivo di \(|m|\) e \(n\). Con il denominatore positivo,
questa forma ridotta è unica; in particolare, lo zero si rappresenta
come \(\frac01\).

In \(\Q\), invece, si assume l'\emph{assioma dell'inverso moltiplicativo}:
\[
  \forall a\in\Q\setminus\{0\}\;\exists b\in\Q
  \quad\text{tale che}\quad
  ab=1.
\]
Con la notazione introdotta in precedenza, tale elemento si indica con
\[
  b=a^{-1}.
\]

L'inverso moltiplicativo di un numero razionale non nullo è detto anche
\emph{reciproco}. Per esempio,
\[
  2^{-1}=\frac12.
\]

\section[Irrazionalità della radice quadrata di 2]{Irrazionalità di \(\sqrt{2}\)}
\label{sec:lez02-irrazionalita-radice-due}

La costruzione dei razionali non è sufficiente a descrivere tutti i numeri
che emergono naturalmente nelle equazioni. Si consideri il polinomio
\[
  p(x)=x^2-2,
  \qquad p\in\Q[x].
\]
La notazione \(\Q[x]\) indica l'insieme dei polinomi nella variabile \(x\)
con coefficienti razionali. L'equazione \(x^2-2=0\) non ammette soluzioni
razionali; in particolare, \(\sqrt{2}\notin\Q\). 

Per affrontare la dimostrazione di questo risultato, si introduce anzitutto la nozione di parità di un intero che 
verrà utilizzata in seguito.

Un intero \(m\) si dice \emph{pari} se esiste \(k\in\Z\) tale che
\[
  m=2k,
\]
mentre si dice \emph{dispari} se esiste \(k\in\Z\) tale che
\[
  m=2k+1.
\]
Ogni intero è pari oppure dispari, e le due possibilità sono mutuamente
esclusive.

\begin{lemma}[Parità del quadrato]
\label{lem:quadrato-pari}
Per ogni \(m\in\Z\), se \(m^2\) è pari, allora \(m\) è pari.
\end{lemma}

\begin{proof}
Richiamando la contronominale introdotta nella
\cref{sec:lez01-insiemi-dimostrazioni}, è sufficiente dimostrare che, se
\(m\) è dispari, allora \(m^2\) è dispari. Poiché \(m\) è dispari, esiste \(k\in\Z\) tale che
\[
  m=2k+1.
\]
Pertanto
\begin{align*}
  m^2
    &=(2k+1)^2\\
    &=4k^2+4k+1\\
    &=4k(k+1)+1.
\end{align*}
Ponendo
\[
  \ell=2k(k+1)\in\Z,
\]
si ottiene
\[
  m^2=2\ell+1,
\]
quindi \(m^2\) è dispari. La contronominale è pertanto verificata e il
lemma segue.
\end{proof}

\begin{theorem}[Irrazionalità di \(\sqrt{2}\)]
\label{thm:irrazionalita-radice-due}
L'equazione \(x^2-2=0\) non ammette soluzioni \(x\in\Q\). Equivalentemente,
\[
  \sqrt{2}\notin\Q.
\]
\end{theorem}

\begin{proof}
Si procede per assurdo, supponendo che l'equazione ammetta una soluzione
razionale. Per la forma ridotta di un numero razionale, esistono
\(m\in\Z\) e \(n\in\N\setminus\{0\}\), con
\[
  \operatorname{mcd}(m,n)=1,
\]
tali che
\[
  \left(\frac{m}{n}\right)^2=2.
\]
Poiché \(n\neq0\), si può moltiplicare per \(n^2\) e ottenere
\[
  m^2=2n^2.
\]
Di conseguenza, \(m^2\) è pari e, per il \cref{lem:quadrato-pari}, anche
\(m\) è pari. Esiste quindi \(h\in\Z\) tale che \(m=2h\). Sostituendo
nell'uguaglianza precedente,
\begin{align*}
  4h^2&=2n^2,\\
  n^2&=2h^2.
\end{align*}
Pertanto anche \(n^2\) è pari e, applicando nuovamente il
\cref{lem:quadrato-pari}, si conclude che \(n\) è pari.

Entrambi gli interi \(m\) e \(n\) sono dunque divisibili per \(2\), in
contraddizione con \(\operatorname{mcd}(m,n)=1\). La contraddizione mostra
che l'ipotesi iniziale è falsa. L'equazione \(x^2-2=0\) non ammette
quindi soluzioni razionali; in particolare,
\[
  \sqrt{2}\notin\Q.
\]

Dal punto di vista logico, la dimostrazione segue lo schema per assurdo
descritto nella \cref{sec:lez01-insiemi-dimostrazioni}. Indicando con
\[
  B:\quad x^2-2=0 \text{ non ammette soluzioni in }\Q,
\]
si assume la negazione \(\neg B\) e da questa si ricava una contraddizione.
Pertanto \(\neg B\) è falsa e la tesi \(B\) è vera.
\end{proof}

Nel risultato appena dimostrato compare il numero \(2\), che è un numero
primo. Prima di formulare la generalizzazione, richiamiamo la relativa
definizione.

\begin{definition}[Numero primo]
\label{def:numero-primo}
Un numero naturale \(p>1\) è primo se i suoi unici divisori positivi sono
\(1\) e \(p\).
\end{definition}

Per esempio, \(2\) e \(3\) sono primi, mentre \(4\) non lo è perché
ammette \(2\) come divisore oltre a \(1\) e a sé stesso. Il numero \(1\)
non è primo, poiché la definizione richiede \(p>1\).

Con questa terminologia, il risultato del
\cref{thm:irrazionalita-radice-due} si generalizza alla radice quadrata
di ogni numero primo.

\begin{proposition}[Irrazionalità della radice di un numero primo]
\label{prop:irrazionalita-radice-primo}
Sia \(p\in\N\) un numero primo. L'equazione
\[
  x^2-p=0
\]
non ammette soluzioni razionali. Equivalentemente,
\[
  \sqrt{p}\notin\Q.
\]
\end{proposition}

La dimostrazione di questa proposizione non viene sviluppata in questa sede.

\section{Fattoriale}
\label{sec:lez02-fattoriale}

Per definire il fattoriale si fissa anzitutto il valore iniziale in
\(n=0\) e si specifica poi come ricavare il valore al naturale successivo.

\begin{definition}[Fattoriale]
\label{def:fattoriale}
Il fattoriale è definito ricorsivamente dalle condizioni
\[
  0!:=1,
  \qquad
  (n+1)!:=(n+1)n!,
  \quad n\in\N.
\]
\end{definition}

Applicando la definizione ricorsiva si ottengono, per esempio,
\begin{align*}
  0!&=1,\\
  1!&=1\cdot0!=1,\\
  2!&=2\cdot1!=2,\\
  3!&=3\cdot2!=3\cdot2\cdot1.
\end{align*}
Per \(n\ge1\), il fattoriale può essere scritto anche nella forma esplicita
\[
  n!=1\cdot2\cdot3\cdots n=\prod_{i=1}^{n}i.
\]
Il caso \(n=0\) è specificato separatamente da \(0!=1\). La ricorrenza
costruisce il prodotto aggiungendo a ogni passo il fattore successivo.

A partire dalla definizione precedente è possibile ricavare alcune
semplici proprietà del fattoriale. In particolare, il fattoriale di un
numero naturale è sempre maggiore o uguale al numero stesso.

\begin{proposition}[Una disuguaglianza per il fattoriale]
\label{prop:fattoriale-maggiore-n}
Per ogni \(n\in\N\) vale
\[
  n!\ge n.
\]
\end{proposition}

\begin{proof}
Per applicare il principio di induzione, introdotto nella
\cref{sec:lez01-induzione}, si considera l'insieme dei numeri naturali
per i quali la disuguaglianza è verificata:
\[
  A=\{n\in\N:n!\ge n\}.
\]

\emph{Passo base.}
Si verifica anzitutto la proprietà per \(n=0\). Per definizione,
\[
  0!=1\ge0,
\]
quindi \(0\in A\) e il passo base è verificato.

\emph{Passo induttivo.}
Si deve dimostrare che
\[
  n\in A\Longrightarrow n+1\in A.
\]

Il caso \(n=0\) viene verificato direttamente:
\[
  1!=1\ge1,
\]
quindi \(1\in A\).

Consideriamo ora \(n\ge1\) e supponiamo \(n\in A\). Per definizione di
\(A\), vale l'ipotesi induttiva
\[
  n!\ge n.
\]
Poiché \(n+1>0\), per la compatibilità dell'ordine con la
moltiplicazione per un numero positivo, espressa nella
\cref{prop:compatibilita-ordine}, moltiplicando entrambi i membri per
\(n+1\) il verso della disuguaglianza non cambia. Si ottiene quindi
\[
  (n+1)n!\ge(n+1)n.
\]
Per la \cref{def:fattoriale},
\[
  (n+1)!=(n+1)n!,
\]
e pertanto
\[
  (n+1)!\ge(n+1)n.
\]

Resta da confrontare \((n+1)n\) con \(n+1\). Poiché \(n\ge1\) e
\(n+1>0\), moltiplicando la disuguaglianza
\[
  n\ge1
\]
per \(n+1\) si ottiene
\[
  (n+1)n\ge(n+1)\cdot1=n+1.
\]
Combinando le due disuguaglianze,
\[
  (n+1)!
  \ge(n+1)n
  \ge n+1.
\]
Si è quindi dimostrato che \(n+1\in A\).

Il passo induttivo è dunque verificato sia per \(n=0\), tramite controllo
diretto, sia per ogni \(n\ge1\).

Essendo verificati il passo base e il passo induttivo, per il
\cref{thm:principio-induzione}, con valore iniziale \(p=0\), si conclude
che
\[
  \N\subseteq A.
\]
Poiché per definizione \(A\subseteq\N\), segue che
\[
  A=\N.
\]
Di conseguenza,
\[
  n!\ge n
\]
per ogni \(n\in\N\).
\end{proof}
